% ═══════════════════════════════════════════════════════════════════════
\section{Qu'est-ce qu'une transformation du plan ?}
% ═══════════════════════════════════════════════════════════════════════

\begin{frame}{Le plan euclidien}
  On travaille dans le \defi{plan euclidien} $\Plan = \R^2$ : un
  ensemble de points, muni d'une distance $d(P,Q)$.
  \[
    P = (x,y) \qquad\text{ou}\qquad P = \col{x}{y} \in \R^2 .
  \]
  \vspace{6pt}
  Tout ce qui suit repose sur \textbf{une seule} définition.
\end{frame}

\begin{frame}{Définition 1.1 --- Transformation géométrique du plan}
  \begin{Enonce}{Définition 1.1 --- Transformation géométrique du plan}
    Une \defi{transformation géométrique} du plan est une
    \textbf{bijection} du plan dans lui-même : une fonction bijective
    $f : \Plan \longrightarrow \Plan$. On note $\Transf$ l'ensemble de
    toutes les transformations du plan.
  \end{Enonce}
  \FigBox[0.42\textheight]{%
    \begin{tikzpicture}[scale=0.95]
      \begin{scope}
        \draw[sectioncolor!50, thick, fill=steelblue4!8]
          (0,0) -- (4,0) -- (5,2.3) -- (1,2.3) -- cycle;
        \node[etiq, sectioncolor,above] at (0.55,2) {$\Plan$};
        \node[pt, label={[etiq]below:$P$}] (P) at (1.4,0.75) {};
        \node[pt, label={[etiq]below:$Q$}] (Q) at (2.6,1.5) {};
      \end{scope}
      \begin{scope}[xshift=6.6cm]
        \draw[firebrick3!50, thick, fill=firebrick3!8]
          (0,0) -- (4,0) -- (5,2.3) -- (1,2.3) -- cycle;
        \node[etiq, firebrick3,above] at (0.55,2) {$\Plan$};
        \node[ptim, label={[etiq]below:$f(P)$}] (Pp) at (2.9,0.6) {};
        \node[ptim, label={[etiq]above:$f(Q)$}] (Qp) at (1.7,1.7) {};
      \end{scope}
      \draw[fleche] (P) to[bend left=12] (Pp);
      \draw[fleche] (Q) to[bend left=12] (Qp);
      \node[etiq, sectioncolor] at (5.6,2.6) {$f$};
    \end{tikzpicture}}
  \begin{center}\small Figure 1.2\end{center}
\end{frame}

\begin{frame}{Décortiquer la définition}
  « Fonction » exige :
  \begin{itemize}
    \item[(i)] $f$ \textbf{partout définie} : chaque point $P$ a une image ;
    \item[(ii)] $f$ \textbf{univoque} : une seule image par point.
  \end{itemize}
  « Bijective » ajoute, \textbf{indépendamment} :
  \begin{itemize}
    \item[(iii)] \textbf{injectivité} : $P \neq Q \Rightarrow f(P) \neq f(Q)$ ;
    \item[(iv)] \textbf{surjectivité} : tout $Q$ possède un antécédent.
  \end{itemize}
  \vspace{6pt}
  \small Quatre conditions indépendantes : les exemples qui suivent
  montrent chacune de leurs violations possibles.
\end{frame}

\begin{frame}{Pour la classe 1.1 --- Trois mises en garde de vocabulaire}
  \begin{Classe}{Trois mises en garde de vocabulaire}
    \begin{enumerate}
      \item Si $Q = f(P)$ : $Q$ est l'\defi{image} de $P$, $P$ un
        \defi{antécédent} de $Q$.
      \item « $P$ s'est déplacé en $Q$ » est une façon de parler
        commode --- mais \textbf{les points ne bougent pas}. Une
        transformation est une \emph{correspondance}, pas un
        mouvement.
      \item Une transformation agit sur \textbf{tout} le plan, pas
        seulement sur la figure dessinée. « Faire tourner un
        triangle », c'est en réalité faire tourner le plan tout
        entier.
    \end{enumerate}
  \end{Classe}
\end{frame}

\begin{frame}{Exemples, et surtout contre-exemples}
  \small Rien ne clarifie mieux une définition que des objets qui n'y
  satisfont pas de justesse.
\end{frame}

\begin{frame}{Exemple 1.1 --- Une symétrie centrale}
  \begin{Enonce}{Exemple 1.1}
    \small Soit $M$ fixe. $f_1(P) = P'$, l'unique point tel que $M$
    soit le milieu de $[PP']$. \textbf{C'est une transformation} :
    chaque point a exactement une image, et chaque point est l'image
    d'exactement un point.
  \end{Enonce}
  \FigBox[0.42\textheight]{%
    \begin{tikzpicture}[scale=1.15]
      \node[pt, label={[etiq]below right:$M$}] (M) at (0,0) {};
      \node[pt, label={[etiq]above right:$P$}] (P) at (1.7,0.9) {};
      \node[ptim, label={[etiq]below left:$P'$}] (Pp) at (-1.7,-0.9) {};
      \draw[sectioncolor!45, thin] (P) -- (Pp);
      \draw[fig] (1.7,0.9) -- (2.3,-0.1) -- (0.9,-0.4) -- cycle;
      \draw[figim] (-1.7,-0.9) -- (-2.3,0.1) -- (-0.9,0.4) -- cycle;
    \end{tikzpicture}}
  \small C'est la \defi{symétrie centrale} de centre $M$ : cas
  particulier d'une homothétie de rapport négatif (section~3).
\end{frame}

\begin{frame}{Exemple 1.2 --- Un repliage du plan}
  \begin{Enonce}{Exemple 1.2}
    \small Soit $d$ fixe. $f_2$ laisse fixe le côté « nord » de $d$, et
    envoie le côté « sud » sur son réfléchi. C'est l'effet de replier
    une feuille sur $d$.
  \end{Enonce}
  \FigBox[0.38\textheight]{%
    \begin{tikzpicture}[scale=1.05]
      \draw[droite] (-2.2,0) -- (2.2,0) node[etiq, right, black] {$d$};
      \foreach \x/\y in {-1.4/1.1, -0.3/0.6, 1.2/0.4}
        {\node[pt] at (\x,\y) {};}
      \foreach \x/\y in {-1.4/-1.1, -0.3/-0.6, 1.2/-0.4}
        {\node[ptim] at (\x,\y) {};
         \draw[gray!60, thin, dotted, -{Stealth[length=1.6mm]}] (\x,\y) -- (\x,-\y);}
    \end{tikzpicture}}
  \textbf{Ce n'est PAS une transformation} : ni injective (deux
  antécédents au nord), ni surjective (rien n'atteint le sud).
\end{frame}

\begin{frame}{Exemple 1.3 --- L'arrondi au pixel près}
  \begin{Enonce}{Exemple 1.3}
    \small $f_3(x,y) = $ le point à coordonnées entières le plus
    proche. C'est ce que fait un écran en affichant une image sur sa
    grille de pixels.
  \end{Enonce}
  \FigBox[0.38\textheight]{%
    \begin{tikzpicture}[scale=0.95]
      \repere{-0.6}{3.6}{-0.6}{2.6}
      \foreach \x/\y in {1/1, 3/2}
        {\fill[firebrick3!18] (\x-0.5,\y-0.5) rectangle (\x+0.5,\y+0.5);
         \node[ptim] at (\x,\y) {};}
    \end{tikzpicture}}
  \textbf{Ce n'est PAS une transformation} : un carré entier s'écrase
  sur un point (pas injective), et les points non entiers ne sont
  jamais atteints (pas surjective).
\end{frame}
